\scanpage{7}
% archive sheet number: 6
\mspage{\handwritten{-3-}}

\uline{Corollaire} (18.9.5) .- \uline{Sous les hypothèses de} (18.9.4) , \uline{le foncteur}

\[
Z\handwritten{\rightsquigarrow} Z\handwritten{\otimes}_A A'\tag*{(18.9.5.1)}
\]

\uline{de la catégorie des préschémas étales sur} $X$ , \uline{dans celle des préschémas étales sur} $X'$ , \uline{est pleinement fidèle} .

Soient $Z_1,Z_2$ deux préschémas étales sur $X$ , et\deleted{\unsure{X}} posons $Z_{\unsure{i|1}}'=Z_i\handwritten{\otimes}_A A'$ $(i=1,2)$ \,; \deleted{\illegible{2}} il s'agit de prouver que tout $X'$-morphisme $Z_1'\handwritten{\longrightarrow} Z_2'$ provient , par changement de base , d'un\inserted{$X$-}morphisme unique $Z_1\handwritten{\longrightarrow} Z_2$ .
\margin{\handwritten{\tikz{\draw (0,0)--(.55,.05) (0,-.11)--(0,.12) (.55,-.06)--(.55,.17);}}}Supposons d'abord que \deleted{$Z_1$ et} $Z_2$ soi\deleted{en}t \uline{séparé\deleted{s}} sur $X$ ; alors, comme $\operatorname{Hom}_X(Z_1,Z_2)$ s'identifie à \deleted{\illegible{1}} $\handwritten{\Gamma}(Z_1\handwritten{\times}_X Z_2/Z_1)$ , et que $Z_1\handwritten{\times}_X\unsure{\mathbb Z|Z}_2$ est étale et séparé sur $Z_1$ , $\operatorname{Hom}_X(Z_1,Z_2)$ s'identifie fonctoriellement, en vertu de (17.9.3) , à l'ensemble des parties ouvertes et fermées $U$ de $Z_1\handwritten{\times}_X Z_2$ telles que la restriction à $U$ de la projection $p:\ Z_1\handwritten{\times}_X Z_2\handwritten{\longrightarrow} Z_1$ soit un morphisme surjectif et radiciel . L'assertion résulte donc dans le cas considéré de (18.9.4) et de ce que $Z_1'\handwritten{\times}_{X'} Z_2'={}= (Z_1\handwritten{\times}_X Z_2)\handwritten{\otimes}_A A'$ .

Passons au cas général : il suffira de prouver que le graphe $\handwritten{\Gamma}'$ d'un $X'$-morphisme $Z_1'\handwritten{\longrightarrow} Z_2'$ \inserted{qui est \uline{ouvert} (17.9.3) \handwritten{dans $Z_1'\times_{X'} Z_2'$}} est de la forme $\handwritten{\Gamma}\handwritten{\otimes}_A A'$ , où $\handwritten{\Gamma}$ est \deleted{\unsure{un ouvert}} induit sur une partie \uline{ouverte} de $Z=Z_1\handwritten{\times}_X Z_2$ . En effet , la restriction à $\handwritten{\Gamma}$ de la projection $p:\ Z\handwritten{\longrightarrow} Z_1$ est alors un \uline{isomorphisme} , car par changement de base \inserted{$A\handwritten{\longrightarrow} A'$} cette restriction devient l'isomorphisme restriction à $\handwritten{\Gamma}'$ de la projection $p':\ Z'\handwritten{\longrightarrow} Z_1'$ (avec $Z'=Z\handwritten{\otimes}_A A'$) , et \margin{\handwritten{\tikz{\draw (0,0) .. controls (.32,.20) and (.50,.03) .. (.54,.21);}}}il suffit d'appliquer (2.7.1,(viii)) ; la conclusion résulte donc \inserted{\handwritten{alors}} de la caractérisation des graphes de morphismes (\uline{I},5.3) .

Si $q:\ Z'\handwritten{\longrightarrow} Z$ est la projection canonique , pour prouver que $\handwritten{\Gamma}'$ est de la forme \deleted{\illegible{1}} $q^{-1}(\handwritten{\Gamma})$ , où $\handwritten{\underline{\Gamma}}$ est ouvert \inserted{dans $Z$} , il suffit , puisque $q$ est un morphisme \inserted{fidèlement} plat \inserted{et quasi-compact} , de prouver \deleted{\unsure{que}} qu'il existe un ensemble $U\handwritten{\subset} Z$ tel que $\handwritten{\Gamma}'=q^{-1}(\unsure{U|V})$ (2.3.12) . Posons $Z''=Z'\handwritten{\times}_Z Z'=Z\handwritten{\times}_X X''$ , où $X''=X'\handwritten{\times}_X X'$ , et soient $q_1:\ Z''\handwritten{\longrightarrow} Z'$ , $q_2:\ Z''\handwritten{\longrightarrow} Z'$ les projections canoniques . Appli-
